\section{Introduction}
\label{sec:introduction}

A developer who adds a dependency expects to download files. What actually
happens in three of the largest language ecosystems is that the client
downloads an archive, unpacks it, and then executes code the package author
supplied, with the developer's full user privileges, before anything has been
imported or called.

This exists for a reason. Packages that wrap native libraries need to compile
against the local toolchain, and install-time hooks are how they do it. The
mechanism is old, load bearing, and enabled by default in every client we
examined. It is also the shortest path from a compromised maintainer account to
code execution on every machine that installs the package, and unlike a
malicious library function it does not require the victim to call anything.

The published incidents are well known and individually
analysed~\citep{ferrante2021incidents,mbeki2022typosquat}. What has not been
measured is the surface itself: how many packages use hooks, what those hooks
actually do when they run, and whether the malicious ones are distinguishable
from the legitimate ones by behaviour rather than by hindsight. Answering that
requires running the hooks, which is why it has mostly not been done. Static
analysis of hook scripts is defeated by two lines of indirection, and every
serious sample we examined used at least that much.

We ran them. \tool{} is an instrumented installer that executes each hook in a
namespaced sandbox with a synthetic home directory, a seeded credential
honeypot, and a network stack that answers every connection while logging it.
Over 19 months we processed every new version published to three registries,
1.42~million versions in total, and recorded 6.1~million hook executions.

\paragraph*{Findings.}
Four results stand out.

\begin{enumerate}[leftmargin=*,itemsep=1pt,topsep=2pt]
  \item Hooks are concentrated where they do the most damage. Only 8.4 percent
        of versions declare one, but those versions account for 31 percent of
        installs by download volume, because hooks live in build tooling that
        sits deep in nearly every dependency tree
        (Section~\ref{sec:results}).
  \item Behaviour separates cleanly. Legitimate compilation hooks touch a
        compiler, the package directory, and nothing else. Only 1.9 percent of
        hooks contact a non-registry network host, and that single feature
        recovers 94 percent of our confirmed malicious set
        (Section~\ref{sec:results}).
  \item The registries hold a strong precursor signal and do not use it. Of
        the 47 packages we assess as malicious, 35 were published within 72
        hours of a maintainer account changing its recovery email address
        (Section~\ref{sec:results}).
  \item Turning hooks off is cheaper than the ecosystem assumes. Installing the
        10{,}000 most downloaded packages with hooks disabled succeeds for 96.2
        percent of them (Section~\ref{sec:defenses}).
\end{enumerate}
